\documentclass[11pt]{article}
\usepackage[utf8]{inputenc}
\usepackage{amsmath,amssymb}
\usepackage[margin=1in]{geometry}
\usepackage{hyperref}
\emergencystretch=3em

\title{Gaussian domination of the fluctuation function:
$S_\lambda(a)\le e^{\beta a^2/2}(\beta/2)^{-\lambda}\Gamma(\lambda)$}
\author{Claude, with Daniel Murfet}
\date{2026-09-06}

\begin{document}
\maketitle
\sloppy

\begin{abstract}
The analytic backbone of grammar \S4.2's tail and convergence estimates, exposed as reusable lemmas:
the complete-the-square domination $e^{-\beta t+\beta a\sqrt t}\le e^{-(\beta/2)t+\beta a^2/2}$ and the
closed-form Gaussian upper bound $S_\lambda(a)\le e^{\beta a^2/2}(\beta/2)^{-\lambda}\Gamma(\lambda)$
for $\lambda>0$. Zero \texttt{sorry}/\texttt{axiom}.
\end{abstract}

\section{Setting}
Every boundary-remainder and convergence bound in \S4.2 (culminating in \texttt{eq:Deltaupper}
$|\Delta|\le Ce^{-\varepsilon n}$) begins with one pointwise inequality: the fluctuation exponent is
sub-Gaussian in $\sqrt t$. Completing the square,
\[
    -\beta t+\beta a\sqrt t = -\tfrac{\beta}{2}t-\tfrac{\beta}{2}(\sqrt t-a)^2+\tfrac{\beta}{2}a^2
      \le -\tfrac{\beta}{2}t+\tfrac{\beta}{2}a^2 ,
\]
so the integrand of $S_\lambda(a)=\int_0^\infty t^{\lambda-1}e^{-\beta t+\beta a\sqrt t}\,dt$ is
dominated by $e^{\beta a^2/2}\,t^{\lambda-1}e^{-(\beta/2)t}$, whose integral is an elementary Gamma
integral. This domination was already buried inside \texttt{fluctuation\_integrableOn}'s proof; here it
is factored out and turned into a quotable bound.

\section{The things formalised}
{\small
\begin{verbatim}
theorem fluctuation_integrand_le (β a t : ℝ) (hβ : 0 < β) (ht : 0 ≤ t) :
    Real.exp (-β*t + β*a*Real.sqrt t) ≤ Real.exp (-(β/2)*t + β*a^2/2)

theorem fluctuation_le_gaussian (β a lam : ℝ) (hβ : 0 < β) (hlam : 0 < lam) :
    fluctuation β lam a ≤ Real.exp (β*a^2/2) * (β/2)^(-lam) * Real.Gamma lam
\end{verbatim}
}

\section{Proof strategy}
The pointwise lemma is \texttt{Real.exp\_le\_exp} composed with \texttt{nlinarith}, fed the two facts
$(\sqrt t)^2=t$ (\texttt{Real.sq\_sqrt}) and $\beta(\sqrt t-a)^2\ge0$ (\texttt{mul\_nonneg} of
\texttt{hβ.le} and \texttt{sq\_nonneg}); \texttt{nlinarith} finds the square completion. The bound then
follows by \texttt{setIntegral\_mono\_on} against the majorant
$e^{\beta a^2/2}\,t^{\lambda-1}e^{-(\beta/2)t^{1}}$ --- integrable by
\texttt{integrableOn\_rpow\_mul\_exp\_neg\_mul\_rpow} at $p=1$ --- and evaluating the majorant integral
by \texttt{integral\_rpow\_mul\_exp\_neg\_mul\_rpow}, which gives
$(\beta/2)^{-\lambda}\Gamma(\lambda)$ at $p=1$, $q=\lambda-1$, $b=\beta/2$.

\section{Roadblocks and resolutions}
\begin{itemize}
\item The Gamma-integral lemmas are stated with the \texttt{rpow} form $t^{1}$ in the exponential; align
the majorant integrand with $t^{(1:\mathbb R)}$ and peel it with \texttt{Real.rpow\_one} inside the
pointwise comparison only.
\item The lemma output carries $/1$ divisions and the exponent $(\lambda-1)+1$; \texttt{simp only
[div\_one]} clears the former and \texttt{rw [show lam - 1 + 1 = lam by ring]} rewrites the exponent in
both the \texttt{rpow} power and the $\Gamma$ argument at once (a bare \texttt{ring} cannot reach under
\texttt{rpow}/$\Gamma$).
\item Splitting $e^{-(\beta/2)t+\beta a^2/2}=e^{\beta a^2/2}e^{-(\beta/2)t}$ needs \texttt{Real.exp\_add}
then \texttt{ring} to reassociate the constant to the front, matching the majorant's shape.
\end{itemize}

\section{What was Mathlib, what was new}
Mathlib supplied the Gamma-integral pair
\texttt{integrableOn\_rpow\_\allowbreak mul\_exp\_\allowbreak neg\_mul\_\allowbreak rpow} and
\texttt{integral\_rpow\_\allowbreak mul\_exp\_\allowbreak neg\_mul\_\allowbreak rpow}, together with
\texttt{setIntegral\_mono\_on}, \texttt{Real.sq\_sqrt}, \texttt{Real.exp\_le\_exp}. New is the exposure of
the domination as a standalone
lemma and the closed-form fluctuation bound --- neither was previously stated, though the majorant
recurred inside the integrability proof.

\section{Lessons}
When an integrability proof already constructs a majorant, its pointwise domination and the value of the
majorant integral are two quotable lemmas waiting to be factored out; do so when downstream estimates
(here the \S4.2 remainder) will lean on the bound repeatedly.

\section{Follow-ups}
The bound is the first ingredient of the boundary-tail remainder $\Delta$ (\texttt{eq:Deltaupper}); the
remainder itself additionally needs the Taylor-tree coefficient bounds (\texttt{prop:CoefficientBound})
and is multivariate.
\end{document}
